\documentclass[12pt, a4paper]{article}
\usepackage{iftex}
\ifPDFTeX
  \usepackage{cmap}
  \usepackage[T2A]{fontenc}
  \usepackage[utf8]{inputenc}
  \usepackage{tempora}
\else
  \usepackage{fontspec}
  \IfFontExistsTF{Times New Roman}
    {\setmainfont{Times New Roman}}
    {\setmainfont{CMU Serif}}
  \IfFontExistsTF{Consolas}
    {\setmonofont{Consolas}[Scale=MatchLowercase]}
    {\IfFontExistsTF{Courier New}
      {\setmonofont{Courier New}}
      {\setmonofont{CMU Typewriter Text}}}
\fi
\usepackage[english,russian]{babel}

\usepackage[a4paper, margin=2.54cm]{geometry}
\setlength{\parindent}{0pt}
\setlength{\parskip}{4pt plus 1pt}
\frenchspacing
\raggedbottom

\usepackage{titlesec}
\setcounter{secnumdepth}{0}
\titleformat{\section}{\large\bfseries}{}{0pt}{}
\titleformat{\subsection}{\normalsize\bfseries}{}{0pt}{}
\titlespacing*{\section}{0pt}{14pt plus 4pt minus 2pt}{4pt plus 2pt}
\titlespacing*{\subsection}{0pt}{10pt plus 2pt minus 2pt}{3pt plus 1pt}

\usepackage[table]{xcolor}
\usepackage{graphicx}
\usepackage{tabularx}
\usepackage{float}
\usepackage{caption}
\DeclareCaptionLabelSeparator{spacedemdash}{~--- }
\captionsetup{labelsep=spacedemdash, font=small, singlelinecheck=false, justification=raggedright}
\captionsetup[table]{position=top, skip=4pt}
\DeclareCaptionType[fileext=lol]{listing}[Листинг][Листинги]
\captionsetup[listing]{position=top, skip=4pt, hypcap=false}
\addto\captionsrussian{\renewcommand{\figurename}{Рисунок}\renewcommand{\tablename}{Таблица}}
\newcommand{\thead}[1]{\multicolumn{1}{|c|}{\cellcolor{gray!20}\textbf{#1}}}
\newcolumntype{L}{>{\raggedright\arraybackslash}X}
\newcolumntype{P}[1]{>{\raggedright\arraybackslash}p{#1}}

\usepackage{fvextra}
\usepackage{needspace}
\fvset{breaklines=true, breakanywhere=true,
       frame=leftline, framerule=1.5pt, framesep=6pt,
       rulecolor=\color{gray!60}, xleftmargin=4pt}
\DefineVerbatimEnvironment{code}{Verbatim}{fontsize=\small}
\DefineVerbatimEnvironment{terminal}{Verbatim}{fontsize=\footnotesize}
\newcommand{\inputterminal}[2][]{\VerbatimInput[fontsize=\footnotesize,#1]{#2}}
\ifPDFTeX
  \usepackage{listingsutf8}
  \lstset{inputencoding=utf8/cp1251, basicstyle=\fontencoding{T2A}\selectfont\ttfamily\small,
          breaklines=true, breakatwhitespace=false, columns=fullflexible,
          keepspaces=true, showstringspaces=false, frame=leftline,
          framerule=1.5pt, framesep=6pt, xleftmargin=4pt,
          rulecolor=\color{gray!60}}
  \newcommand{\inputcode}[2][]{\lstinputlisting[#1]{#2}}
  \newcommand{\inputappendixcode}[2][]{\lstinputlisting[basicstyle=\linespread{0.90}\fontencoding{T2A}\selectfont\ttfamily\footnotesize,#1]{#2}}
\else
  \newcommand{\inputcode}[2][]{\VerbatimInput[fontsize=\small,#1]{#2}}
  \newcommand{\inputappendixcode}[2][]{\VerbatimInput[fontsize=\footnotesize,baselinestretch=0.90,#1]{#2}}
\fi
\newenvironment{listingblock}[1]
  {\Needspace{8\baselineskip}\par\medskip\captionof{listing}{#1}\nopagebreak\vspace{-\topsep}\vspace{-\parskip}}
  {\par}

\newcounter{appendixno}
\newcommand{\appendixsection}[2]{\clearpage\section{Приложение #1. #2}\stepcounter{appendixno}\setcounter{listing}{0}\renewcommand{\thelisting}{#1.\arabic{listing}}\renewcommand{\theHlisting}{app\arabic{appendixno}.\arabic{listing}}}

\usepackage[hidelinks]{hyperref}
\usepackage{xurl}

\newcommand{\hint}[1]{\leavevmode{\color{gray}\itshape #1}}
\newcommand{\placeholder}[2][4.5cm]{\fbox{\parbox[c][#1][c]{0.95\linewidth}{\centering\hint{#2}}}}


% ===================== Данные для титульного листа =====================
\newcommand{\ReportLab}{2}
\newcommand{\ReportStudent}{Новак Дмитрий Евгеньевич}
\newcommand{\ReportGroup}{М8О-201БВ-25}
\newcommand{\ReportTeacher}{Миронов Евгений Сергеевич}
\newcommand{\ReportDate}{06.10.2026}
\newcommand{\ReportYear}{2026}
\newcommand{\ReportCourse}{Операционные системы и архитектура компьютера}
% ========================================================================

\begin{document}
% Титульный лист. Данные берутся из команд \Report... в report.tex.
\begin{titlepage}
  \centering\large

  МОСКОВСКИЙ АВИАЦИОННЫЙ ИНСТИТУТ\\
  (НАЦИОНАЛЬНЫЙ ИССЛЕДОВАТЕЛЬСКИЙ УНИВЕРСИТЕТ)\par
  \vspace{18pt}
  Институт №8 «Компьютерные науки и прикладная математика»\\
  Кафедра 806 «Вычислительная математика и программирование»\par

  \vspace{\stretch{1.2}}

  {\bfseries Лабораторная работа №\,\ReportLab{} по курсу\\
  «\ReportCourse»\par}

  \vspace{\stretch{1.6}}

  \begin{flushright}
    \begin{tabular}{r@{\hspace{0.6em}}l}
      Выполнил:      & \ReportStudent \\[2pt]
      Группа:        & \ReportGroup   \\[2pt]
      Преподаватель: & \ReportTeacher \\[2pt]
      Дата:          & \ReportDate    \\[2pt]
      Оценка:        & \rule[-2pt]{4.2cm}{0.4pt} \\
    \end{tabular}
  \end{flushright}

  \vspace{\stretch{1}}

  Москва, \ReportYear
\end{titlepage}
\setcounter{page}{2}


% ------------------------------------------------------------------------
\section{Условие}

\textbf{Цель работы:} изучить управление процессами в ОС семейства UNIX и
организацию межпроцессного взаимодействия с помощью неименованных каналов (pipe),
а также получить практические навыки работы с системными вызовами создания,
запуска и ожидания процессов, перенаправления стандартных потоков и обработки
системных ошибок.

\textbf{Задание:} родительский процесс создаёт дочерний процесс. Пользователь
задаёт путь к входному файлу; файл открывается на чтение и становится стандартным
вводом дочернего процесса. Стандартный вывод дочернего процесса направляется в
канал. Родительский процесс читает данные из канала и выводит их в свой стандартный
поток вывода. Родительская и дочерняя части реализуются отдельными программами.
Ошибки системных вызовов должны обрабатываться.

\textbf{Вариант 9:} каждая строка входного файла содержит произвольное количество
чисел типа \texttt{float}. Дочерний процесс делит первое число строки последовательно
на все следующие числа и выводит результат. Проверка деления на ноль выполняется
в дочернем процессе; при обнаружении нулевого делителя выполнение завершается с
ошибкой.

% ------------------------------------------------------------------------
\section{Метод решения}

Программа состоит из двух исполняемых файлов: \texttt{lab2\_parent} и
\texttt{lab2\_child}. Родитель сначала считывает введённый пользователем путь и
открывает соответствующий файл только для чтения. Затем создаётся неименованный
канал и дочерний процесс.

После \texttt{fork()} дочерний процесс перенаправляет открытый файл на свой
\texttt{stdin} с помощью \texttt{dup2()}, а записывающий конец канала --- на свой
\texttt{stdout}. После закрытия лишних дескрипторов выполняется
\texttt{execl()}, заменяющий код дочернего процесса программой
\texttt{lab2\_child}. Дочерняя программа читает стандартный ввод и пишет в
стандартный вывод. Источник ввода и получатель вывода определяются перенаправлением
дескрипторов.

Дочерняя программа читает файл построчно функцией \texttt{getline()}, преобразует
числа с помощью \texttt{strtof()} и последовательно выполняет деление. Нулевой
делитель проверяется до операции деления. Результат каждой строки выводится в
\texttt{stdout}, поэтому фактически попадает в канал.

Родитель закрывает неиспользуемый записывающий конец канала и в цикле читает
результаты из читающего конца. Полученные байты записываются в стандартный вывод
родителя. После получения EOF из канала родитель закрывает дескриптор и ожидает
завершения дочернего процесса с помощью \texttt{waitpid()}, после чего анализирует
код его завершения.

% ------------------------------------------------------------------------
\section{Описание программы}

\subsection{Структура программной части}

\begin{table}[H]
  \caption{Файлы программной части}
  \begin{tabularx}{\linewidth}{|P{4.2cm}|L|}
    \hline
    \thead{Файл} & \thead{Назначение} \\ \hline
    \texttt{src/parent.c} & Родительский процесс: получает путь к файлу, открывает его, создаёт канал и дочерний процесс, получает результаты из канала и ожидает завершения ребёнка. \\ \hline
    \texttt{src/child.c} & Дочерняя программа: читает строки из стандартного ввода, выполняет деление первого числа на последующие и выводит результат. \\ \hline
    \texttt{CMakeLists.txt} & Настройка сборки двух исполняемых файлов \texttt{lab2\_parent} и \texttt{lab2\_child}. \\ \hline
    \texttt{data/*.txt} & Тестовые входные файлы для корректного ввода, отрицательных чисел и деления на ноль. \\ \hline
  \end{tabularx}
\end{table}

Исходный код в листингах подключается напрямую из \texttt{src/}, без копий в
каталоге отчёта. Трассировка разобрана в отдельном разделе и приложении Б.

\subsection{Основные типы данных и функции}

В родительском процессе путь к входному файлу считывается функцией
\texttt{getline()} в динамически выделяемый буфер \texttt{filename}. После
успешного \texttt{open()} строка пути больше не требуется и освобождается функцией
\texttt{free()}. Дескриптор входного файла хранится в \texttt{input\_fd}, а два
конца канала --- в массиве \texttt{pipe\_fd[2]}.

Для передачи данных родитель использует буфер \texttt{char buffer[4096]}.
Дочерний процесс хранит текущую строку в динамическом буфере \texttt{line}, который
переиспользуется между вызовами \texttt{getline()} и освобождается перед выходом.

В \texttt{child.c} вспомогательная функция
\texttt{read\_float(char* text, char** end, float* value)} преобразует очередное
число через \texttt{strtof()}, проверяет отсутствие преобразования, ошибку
\texttt{ERANGE} и нечисловые/бесконечные значения. Возвращаемое значение 1 означает
успешное чтение числа, 0 --- ошибку.

Программа использует отдельные коды завершения: 1 --- общая ошибка, 2 --- деление
на ноль, 3 --- некорректные входные данные, 4 --- ошибка запуска дочерней программы
через \texttt{exec}.

\subsection{Используемые системные вызовы}

\begin{table}[H]
  \caption{Используемые системные вызовы и POSIX-интерфейсы}
  \begin{tabularx}{\linewidth}{|P{5.2cm}|L|}
    \hline
    \thead{Вызов} & \thead{Назначение в программе} \\ \hline
    \texttt{open(path, O\_RDONLY)} & Открытие указанного пользователем входного файла только для чтения. \\ \hline
    \texttt{pipe(pipe\_fd)} & Создание однонаправленного канала между дочерним и родительским процессами. \\ \hline
    \texttt{fork()} & Создание дочернего процесса. \\ \hline
    \texttt{dup2(oldfd, newfd)} & Перенаправление стандартного ввода ребёнка на файл и стандартного вывода ребёнка в канал. \\ \hline
    \texttt{execl()} & Замена кода дочернего процесса программой \texttt{lab2\_child}. \\ \hline
    \texttt{read()} & Чтение родителем результатов из канала. \\ \hline
    \texttt{write()} & Вывод родителем полученных из канала байтов в \texttt{stdout}. \\ \hline
    \texttt{close()} & Закрытие неиспользуемых и более не нужных файловых дескрипторов. \\ \hline
    \texttt{waitpid()} & Ожидание завершения конкретного дочернего процесса и получение его статуса. \\ \hline
  \end{tabularx}
\end{table}

\subsection{Ключевые фрагменты кода}

Сначала родитель открывает заданный пользователем файл, создаёт канал и дочерний
процесс. Ошибки каждого системного вызова проверяются.

\begin{listingblock}{Открытие файла, создание канала и процесса (parent.c)}
  \inputcode[firstline=48, lastline=88]{../src/parent.c}
\end{listingblock}

В дочерней ветви закрывается читающий конец канала. Затем входной файл дублируется
в дескриптор 0, а записывающий конец канала --- в дескриптор 1. После этого
\texttt{execl()} запускает отдельную дочернюю программу.

\begin{listingblock}{Перенаправление стандартных потоков и запуск child (parent.c)}
  \inputcode[firstline=90, lastline=120]{../src/parent.c}
\end{listingblock}

Родитель читает поток байтов из канала до EOF и передаёт полученные данные в свой
стандартный вывод.

\begin{listingblock}{Чтение результатов из канала (parent.c)}
  \inputcode[firstline=122, lastline=177]{../src/parent.c}
\end{listingblock}

Дочерняя программа проверяет нулевой делитель до выполнения операции деления.

\begin{listingblock}{Проверка деления на ноль (child.c)}
  \inputcode[firstline=64, lastline=86]{../src/child.c}
\end{listingblock}

% ------------------------------------------------------------------------
\section{Результаты}

Проект собирался CMake с включёнными предупреждениями \texttt{-Wall -Wextra}.
После сборки оба исполняемых файла находятся в одной директории
\texttt{build/lab\_02}. Команды запуска ниже выполняются из этой директории:
путь к входному файлу пользователь вводит сам, а \texttt{./lab2\_child} ищется
относительно текущей рабочей директории.

Пример успешного запуска на тестовом файле \texttt{ok.txt}. После приглашения
«Введите путь к файлу:» пользователь вводит указанный путь:

\begin{terminal}
$ ./lab2_parent
../../lab_02/data/ok.txt
0.5
20
10
100000
500
0.05
0
\end{terminal}

Дополнительно были проверены файлы \texttt{negatives.txt} с отрицательными
числами и \texttt{zero.txt} с нулевым делителем. Отрицательные значения
корректно обрабатываются как значения типа
\texttt{float}. При появлении нуля среди делителей дочерний процесс выводит
сообщение об ошибке в \texttt{stderr} и завершается с кодом 2; родитель получает
этот статус через \texttt{waitpid()} и завершает работу с тем же кодом. При
проверочном запуске \texttt{zero.txt} сообщение в \texttt{stderr} имеет вид
\texttt{Ошибка в строке 1: деление на ноль.}

Проверено перенаправление стандартных потоков оболочкой. Следующая команда
скрывает обычный вывод и оставляет видимыми только сообщения об ошибках:

\begin{terminal}
$ echo "../../lab_02/data/zero.txt" | ./lab2_parent > /dev/null
\end{terminal}

А команда ниже, наоборот, скрывает \texttt{stderr} и оставляет \texttt{stdout}:

\begin{terminal}
$ echo "../../lab_02/data/zero.txt" | ./lab2_parent 2> /dev/null
\end{terminal}

Потоки также можно сохранить в разные файлы:

\begin{terminal}
$ echo "../../lab_02/data/zero.txt" | ./lab2_parent > results.txt 2> errors.txt
\end{terminal}

При таком запуске с
\texttt{zero.txt} код завершения равен 2, в \texttt{errors.txt} находится
сообщение о делении на ноль, а в \texttt{results.txt} --- только приглашение
ввести путь.

% ------------------------------------------------------------------------
\section{Демонстрация работы strace}

Трассировка выполнена в Ubuntu ARM64 командой \texttt{strace -f}; параметр
\texttt{-f} необходим, чтобы продолжить трассировку после \texttt{fork()} и видеть
как родительский, так и дочерний процесс. Ниже приведены фрагменты файла
\path{trace/strace_short.txt}. Полная трассировка аналогичного запуска
\path{trace/strace.txt} приведена в приложении Б; она получена отдельно,
поэтому PID процессов в двух файлах различаются.

Сначала родитель считывает путь, открывает файл, создаёт канал и дочерний процесс:

\inputterminal[firstline=8,lastline=11]{trace/strace_short.txt}

В исходном коде используются \texttt{open()}, \texttt{pipe()} и \texttt{fork()},
однако \texttt{strace} показывает фактические системные вызовы Linux. В данной
системе соответствующие обёртки libc реализованы через \texttt{openat()},
\texttt{pipe2()} и \texttt{clone()}. В результате файл получил дескриптор 3,
читающий конец канала --- 4, записывающий --- 5; PID ребёнка равен 320722.

Далее дочерний процесс перенаправляет дескрипторы и запускает
\texttt{lab2\_child}:

\inputterminal[firstline=17,lastline=23]{trace/strace_short.txt}

В программе вызван \texttt{dup2()}, но на Ubuntu ARM64 libc использует системный
вызов \texttt{dup3()} с нулевыми флагами. После первого вызова дескриптор 0
(\texttt{stdin}) указывает на входной файл, после второго дескриптор 1
(\texttt{stdout}) --- на записывающий конец канала. Вызов \texttt{execl()} в
исходном коде приводит к системному вызову \texttt{execve()} и не создаёт новый
процесс: PID 320722 сохраняется.

Передача результата через канал видна непосредственно в трассировке:

\inputterminal[firstline=30,lastline=39]{trace/strace_short.txt}

Дочерний процесс пишет результаты в свой \texttt{stdout}, который направлен в
канал. Родитель читает те же байты из дескриптора 4 и выводит их на экран. При этом
несколько отдельных операций \texttt{write()} ребёнка могут быть получены одним
\texttt{read()} родителя, поскольку pipe представляет собой поток байтов, а не
набор сообщений.

После завершения ребёнка родитель получает EOF из канала, закрывает его и ожидает
дочерний процесс:

\inputterminal[firstline=48,lastline=52]{trace/strace_short.txt}

В исходном коде используется \texttt{waitpid()}, а трассировка Linux показывает
\texttt{wait4()}. Ребёнок завершился с кодом 0; затем с тем же кодом завершился
родитель.

% ------------------------------------------------------------------------
\section{Выводы}

Таким образом, данная лабораторная работа никак кардинально не
изменила мою жизнь, но позволила более подробно изучить работу процессов и
их взаимодействие через каналы, а также узнать, что на самом деле даже самые
большие программы являются fork-ами от чего-то другого.

% ========================================================================
\appendixsection{А}{Исходный код программы}

\begin{listingblock}{Файл parent.c}
  \inputappendixcode{../src/parent.c}
\end{listingblock}

\begin{listingblock}{Файл child.c}
  \inputappendixcode{../src/child.c}
\end{listingblock}

\begin{listingblock}{Файл CMakeLists.txt}
  \inputappendixcode{../CMakeLists.txt}
\end{listingblock}

% ========================================================================
\appendixsection{Б}{Полный вывод strace}

Полный вывод получен при запуске программы на Ubuntu ARM64 с параметром
\texttt{-f}, позволяющим трассировать дочерний процесс.

\begin{listingblock}{Полный вывод strace -f ./lab2\_parent}
  \inputterminal[baselinestretch=0.95]{trace/strace.txt}
\end{listingblock}

\end{document}
